% Bagian: computability
% Bab: recursive-functions
% Subbagian: pr-functions-computable

\documentclass[../../../include/open-logic-section]{subfiles}

\begin{document}

\olfileid[id]{cmp}{rec}{cmp}
\olsection{Fungsi Rekursif Primitif Dapat Dihitung}

Andaikan suatu fungsi $h$ didefinisikan dengan rekursi primitif
\begin{eqnarray*}
h(\vec x, 0)     & = & f(\vec x) \\
h(\vec x, y + 1) & = & g(\vec x, y, h(\vec x, y))
\end{eqnarray*}
dan andaikan fungsi $f$ serta $g$ dapat dihitung. (Kita menggunakan $\vec x$
sebagai singkatan untuk $x_0$, \dots, $x_{k-1}$.) Maka $h(\vec x, 0)$ jelas
dapat dihitung karena nilainya hanyalah $f(\vec x)$, yang kita andaikan dapat
dihitung. Selanjutnya, $h(\vec x, 1)$ juga dapat dihitung karena $1 = 0 + 1$,
sehingga $h(\vec x, 1)$ hanyalah
\begin{align*}
 h(\vec x, 1) & = g(\vec x, 0, h(\vec x, 0)) =  g(\vec x, 0, f(\vec x)).
\intertext{Kita dapat melanjutkan dengan cara ini dan menghitung}
h(\vec x, 2) & = g(\vec x, 1, h(\vec x, 1)) = g(\vec x, 1, g(\vec x, 0, f(\vec x)))\\
h(\vec x, 3) & = g(\vec x, 2, h(\vec x, 2)) = g(\vec x, 2, g(\vec x, 1, g(\vec x, 0, f(\vec x))))\\
h(\vec x, 4) & = g(\vec x, 3, h(\vec x, 3)) = g(\vec x, 3, g(\vec x, 2, g(\vec x, 1, g(\vec x, 0, f(\vec x)))))\\
& \vdots
\end{align*}
Jadi, untuk menghitung $h(\vec x, y)$ secara umum, hitunglah secara
berturut-turut $h(\vec x, 0)$, $h(\vec x, 1)$, \dots, hingga kita mencapai
$h(\vec x, y)$.

Dengan demikian, suatu definisi rekursif primitif menghasilkan fungsi dapat
dihitung yang baru jika fungsi $f$ dan $g$ dapat dihitung. Komposisi fungsi
juga menghasilkan fungsi dapat dihitung jika fungsi $f$ dan $g_i$ dapat
dihitung.

Karena fungsi dasar $\Zero$, $\Succ$, dan $\Proj{n}{i}$ dapat dihitung,
sedangkan komposisi dan rekursi primitif menghasilkan fungsi dapat dihitung
dari fungsi-fungsi dapat dihitung, setiap fungsi rekursif primitif dapat
dihitung.

\end{document}

